\section{System Model}
\label{sec:system}

\defname{} instantiates an orchestrated collective of LLM agents. Because it is an
environment, the model below is \emph{configurable}: the quantities that a security
study wants to vary---topology, roles, observability, trust---are parameters of the
harness, not fixed choices.

\subsection{Agents and roles}
The system comprises $N\in\{3,\dots,50\}$ agents $\Agents=\{a_1,\dots,a_N\}$, each an
LLM policy that reads a local context and emits messages and tool calls. Agents are
\emph{heterogeneous} in role, drawn from $\Roles$:
\begin{itemize}[leftmargin=1.4em]
\item an \textbf{orchestrator} (planner) that decomposes the task, assigns subtasks,
and routes messages;
\item \textbf{workers} that solve subtasks and exchange intermediate results;
\item \textbf{tool-executors} that mediate calls to external tools; and
\item a \textbf{shared-memory store} that agents read from and write to.
\end{itemize}
Each agent $a$ has capabilities: a model backbone (\Cref{sec:experiments}), a set of
permitted tools, and read/write permissions on shared memory. The role assignment and
capability grants are part of the configuration $\config$.

\subsection{Topology and communication}
Agents communicate over a directed graph $\Gtop=(\Vtop,\Etop)$ with $\Vtop=\Agents$;
an edge $(a,a')\in\Etop$ means $a'$ may read messages from $a$. \defname{} ships four
canonical topologies (\Cref{fig:architecture}): a \textbf{star} with the orchestrator at
the hub; a \textbf{chain} in which each agent forwards to its successor; a \textbf{tree}
in which the orchestrator branches to sub-orchestrators and workers; and a \textbf{mesh}
in which many worker pairs communicate directly. Topology is a first-class variable
because, as \Cref{sec:theory} shows, it controls how a compromise propagates and how
severe an orchestrator compromise is.

\subsection{Partial observability}
Agents are \emph{partially observing}: agent $a$ sees only its local context---its own
role prompt, the messages on its inbound edges $\Nin(a)$, and the memory entries it is
permitted to read. No agent observes the global state $s_t$; in particular, no honest
agent can unilaterally detect that a peer has been compromised or that the orchestrator
is malicious. Partial observability is what makes group-level detection nontrivial and
is a stated requirement of the multi-agent security agenda~\cite{schroederdewitt2025masec}.

\subsection{Forced coordination}
A task is \emph{forced-coordination} when its benign-success predicate $\Predutil$ can be
satisfied only by the joint action of several agents---no single agent, and no coalition
below a task-specific size, can complete it alone (\Cref{def:forced} makes this precise).
Forced coordination is the property that makes the collective the unit of analysis: it
cannot be represented in a single-agent harness, and it lets \defname{} measure
coordination quality and coordination sabotage as system-level quantities.

\subsection{Orchestration, trust boundaries, and the harness}
An \textbf{orchestration layer}---part of the trusted harness---instantiates $\Gtop$,
assigns roles, enforces the observability masks, applies the adversary knobs, and drives
episodes step by step. Per-agent \textbf{trust boundaries} record which agents and tools
a given agent is configured to trust; the orchestrator itself carries a trust bit that
can be set to \emph{trusted} or \emph{compromised} (\Cref{sec:threat}). Optional
\textbf{human-approval hooks} can gate designated high-risk actions, letting a study
measure the human-intervention rate required to hold risk at a target level.

\subsection{Trusted computing base}
\label{sec:tcb}
Following the source design, the \textbf{deterministic state checker} $\Checker$ and the
\textbf{environment harness} (orchestration layer, logging, observability enforcement)
are the trusted computing base (TCB) and are assumed uncompromised; the emulated tools
have \emph{no real side effects}. Everything else---agents (including the orchestrator),
tools, memory contents, and messages---may be adversarial under a given configuration.
Out of scope are real-world tool side effects, physical actuation, and model-weight
attacks. This TCB is what makes hijack-resistant scoring possible: because $\Checker$
reads only the harness-recorded trace $\Trace$, no adversarial agent or tool can alter a
verdict except by actually changing the security-relevant state that the harness records
(\Cref{thm:integrity}).
